\documentclass[10pt,a4paper]{amsart}
\usepackage{amsmath,amssymb,amsthm,graphicx}
\usepackage{verbatim}
\newenvironment{reference}{\comment}{\endcomment}
%\newenvironment{reference}{}{}
\newenvironment{slogan}{\comment}{\endcomment}
\newenvironment{history}{\comment}{\endcomment}

% For commutative diagrams we use Xy-pic
\usepackage[all]{xy}

% We use 2cell for 2-commutative diagrams.
\xyoption{2cell}
% Theorem environments.
%
\theoremstyle{plain}
\newtheorem{theorem}[subsection]{Theorem}
\newtheorem{proposition}[subsection]{Proposition}
\newtheorem{lemma}[subsection]{Lemma}

\theoremstyle{definition}
\newtheorem{definition}[subsection]{Definition}
\newtheorem{example}[subsection]{Example}
\newtheorem{exercise}[subsection]{Exercise}
\newtheorem{situation}[subsection]{Situation}

\theoremstyle{remark}
\newtheorem{remark}[subsection]{Remark}
\newtheorem{remarks}[subsection]{Remarks}

\numberwithin{equation}{subsection}

% Macros
%
\def\lim{\mathop{\mathrm{lim}}\nolimits}
\def\colim{\mathop{\mathrm{colim}}\nolimits}
\def\Spec{\mathop{\mathrm{Spec}}}
\def\Hom{\mathop{\mathrm{Hom}}\nolimits}
\def\Ext{\mathop{\mathrm{Ext}}\nolimits}
\def\SheafHom{\mathop{\mathcal{H}\!\mathit{om}}\nolimits}
\def\SheafExt{\mathop{\mathcal{E}\!\mathit{xt}}\nolimits}
\def\Sch{\mathit{Sch}}
\def\Mor{\mathop{\mathrm{Mor}}\nolimits}
\def\Ob{\mathop{\mathrm{Ob}}\nolimits}
\def\Sh{\mathop{\mathit{Sh}}\nolimits}
\def\NL{\mathop{N\!L}\nolimits}
\def\CH{\mathop{\mathrm{CH}}\nolimits}
\def\proetale{{pro\text{-}\acute{e}tale}}
\def\etale{{\acute{e}tale}}
\def\QCoh{\mathit{QCoh}}
\def\Ker{\mathop{\mathrm{Ker}}}
\def\Im{\mathop{\mathrm{Im}}}
\def\Coker{\mathop{\mathrm{Coker}}}
\def\Coim{\mathop{\mathrm{Coim}}}

% Boxtimes
%
\DeclareMathSymbol{\boxtimes}{\mathbin}{AMSa}{"02}

%
% Macros for moduli stacks/spaces
%
\def\QCohstack{\mathcal{QC}\!\mathit{oh}}
\def\Cohstack{\mathcal{C}\!\mathit{oh}}
\def\Spacesstack{\mathcal{S}\!\mathit{paces}}
\def\Quotfunctor{\mathrm{Quot}}
\def\Hilbfunctor{\mathrm{Hilb}}
\def\Curvesstack{\mathcal{C}\!\mathit{urves}}
\def\Polarizedstack{\mathcal{P}\!\mathit{olarized}}
\def\Complexesstack{\mathcal{C}\!\mathit{omplexes}}
% \Pic is the operator that assigns to X its picard group, usage \Pic(X)
% \Picardstack_{X/B} denotes the Picard stack of X over B
% \Picardfunctor_{X/B} denotes the Picard functor of X over B
\def\Pic{\mathop{\mathrm{Pic}}\nolimits}
\def\Picardstack{\mathcal{P}\!\mathit{ic}}
\def\Picardfunctor{\mathrm{Pic}}
\def\Deformationcategory{\mathcal{D}\!\mathit{ef}}
\usepackage{cleveref}
\crefname{originaldefinition}{original definition}{original definitions}
\Crefname{originaldefinition}{Original definition}{Original definitions}
\makeatletter
\newlabel{algebra-definition-complete}{{0317}{source}{}{sourcepointer-1}{}}
\newlabel{algebra-definition-complete@cref}{{[originaldefinition][0][]0317}{[1][0][]source}{}{}{}}
\makeatother
\crefname{originallemma}{original lemma}{original lemmas}
\Crefname{originallemma}{Original lemma}{Original lemmas}
\makeatletter
\newlabel{algebra-lemma-hathat-finitely-generated}{{05GG}{source}{}{sourcepointer-2}{}}
\newlabel{algebra-lemma-hathat-finitely-generated@cref}{{[originallemma][0][]05GG}{[1][0][]source}{}{}{}}
\makeatother
\crefname{originalreference}{original reference}{original references}
\Crefname{originalreference}{Original reference}{Original references}
\makeatletter
\newlabel{lemma-closed-is-complete}{{0HB3}{source}{}{sourcepointer-3}{}}
\newlabel{lemma-closed-is-complete@cref}{{[originalreference][0][]0HB3}{[1][0][]source}{}{}{}}
\makeatother
\title{The category of I-adically complete modules over a Noetherian ring has kernels and cokernels but is not abelian}
\author{The Stacks Project Authors}
\date{stacks-project:examples.tex:local-snapshot-20261012; GNU FDL 1.2 or later}
\begin{document}
\maketitle
\noindent\textit{Editorial note.} All delivered mathematical text is the author's own, copied from the pinned Stacks Project chapter file examples.tex; no statement or proof step was rewritten, reordered or removed.\par\medskip
\noindent\textit{Editorial note.} The counted proof is the author's complete argument: The closure argument uses the p-adic topology on the product; the delivered card L48284 (closed-is-complete) is the chapter's standing tool for closedness of kernels in complete modules.\par\medskip
\noindent\textit{Editorial note.} The label lemma-complete-modules-not-abelian names the counted statement, and the proof environment immediately following it in the source is delivered verbatim as its proof.\par\medskip
\noindent\textit{Editorial note.} The counted statement is carried by the author's own lemma environment in the chapter source.\par\medskip
\noindent\textit{Allowed standard prerequisites (editorial).}\begin{itemize}
\item Context span 792-805 (section-non-abelian preamble, shared with delivered L48308): for I finitely generated the category A of I-adically complete R-modules has cokernel of phi: M -> N given by the I-adic completion of the usual cokernel (Algebra, hathat-finitely-generated) and kernel the usual kernel, which is closed and I-adically complete by Lemma closed-is-complete (delivered card L48284)
\item Context span 841-859: with R = Z\_p = lim Z/p\textasciicircum{}n Z and I = (p), the map diag(1, p, p\textasciicircum{}2, ...) from (direct sum of Z\_p)\textasciicircum{}wedge to the product of Z\_p sends (x\_1, x\_2, ...) with v\_p(x\_i) -> infinity to (x\_1, p x\_2, p\textasciicircum{}2 x\_3, ...); the element (1, p, p\textasciicircum{}2, ...) lies in the closure of the image but not in the image; by the kernel/cokernel description this shows Im != Coim for this map
\end{itemize}\medskip
\ifdefined\hypertarget\hypertarget{sourcepointer-1}{}\fi
\noindent\textit{Source reference.} Original definition 0317: sources/stacks-chow/tags.txt, source lines 3974--3974. Prerequisite pointer invoked at 875; it names a standard fact of a sibling chapter that the counted argument uses, and every step of the counted proof built on it is retained verbatim in the delivered body.\par\medskip
\ifdefined\hypertarget\hypertarget{sourcepointer-2}{}\fi
\noindent\textit{Source reference.} Original lemma 05GG: sources/stacks-chow/tags.txt, source lines 7060--7060. Prerequisite pointer invoked at 880; it names a standard fact of a sibling chapter that the counted argument uses, and every step of the counted proof built on it is retained verbatim in the delivered body.\par\medskip
\ifdefined\hypertarget\hypertarget{sourcepointer-3}{}\fi
\noindent\textit{Source reference.} Original reference 0HB3: sources/stacks-examples/pinned-original-source.tex, source lines 808--808. Prerequisite pointer invoked at 882; it names a standard fact of the same chapter that the counted argument uses, and every step of the counted proof built on it is retained verbatim in the delivered body.\par\medskip
\noindent
Let $R$ be a ring and let $I \subset R$ be a finitely generated ideal.
Consider the category $\mathcal{A}$ of $I$-adically complete
$R$-modules, see
Algebra, Definition \ref{algebra-definition-complete}.
Let $\varphi : M \to N$ be a morphism of $\mathcal{A}$.
The cokernel of $\varphi$ in $\mathcal{A}$ is the $I$-adic completion
$(\Coker(\varphi))^\wedge$ of the usual cokernel
(as $I$ is finitely generated this completion is complete, see
Algebra, Lemma \ref{algebra-lemma-hathat-finitely-generated}).
Let $K = \Ker(\varphi)$; since $\varphi$ is continuous this
is a closed submodule of $M$. By Lemma \ref{lemma-closed-is-complete}
below we see that $K$ is $I$-adically complete and hence is the
kernel of $\varphi$ in $\mathcal{A}$.

\noindent
We will give an example that shows that $\Im \not = \Coim$ in the
category $\mathcal{A}$ in general. We take
$R = \mathbf{Z}_p = \lim_n \mathbf{Z}/p^n\mathbf{Z}$
to be the ring of $p$-adic integers and we take $I = (p)$.
Consider the map
$$
\text{diag}(1, p, p^2, \ldots) :
\left(\bigoplus\nolimits_{n \geq 1} \mathbf{Z}_p\right)^\wedge
\longrightarrow
\prod\nolimits_{n \geq 1} \mathbf{Z}_p
$$
where the left hand side is the $p$-adic completion of the direct sum.
Hence an element of the left hand side is a vector $(x_1, x_2, x_3, \ldots)$
with $x_i \in \mathbf{Z}_p$ with $p$-adic valuation $v_p(x_i) \to \infty$ as
$i \to \infty$. This maps to $(x_1, px_2, p^2x_3, \ldots)$. Hence we see
that $(1, p, p^2, \ldots)$ is in the closure of the image but not in
the image. By our description of kernels and cokernels above it is
clear that $\Im \not = \Coim$ for this map.

\begin{lemma}
\label{lemma-complete-modules-not-abelian}
Let $R$ be a ring and let $I \subset R$ be a finitely generated ideal.
The category of $I$-adically complete $R$-modules has kernels and
cokernels but is not abelian in general even when $R$ is Noetherian.
\end{lemma}

\begin{proof}
See above.
\end{proof}

\end{document}
